\section{Conclusion}
\label{sec:conclusion}

We have presented a depth-aware polka dot art generation pipeline that
integrates YOLO26 monocular depth estimation and semantic segmentation
with Weighted Voronoi Stippling to produce Yayoi Kusama--style
non-photorealistic renderings from single RGB images. The central
contribution is a mathematically principled depth-to-radius mapping
function that encodes scene depth into dot size---near objects receive
large, prominent dots while far objects recede into fields of small,
sparse dots---with controlled Gaussian randomization that replicates
the organic variation of hand-painted dots. This mapping is governed
by two intuitive parameters: a depth falloff exponent $\alpha$ that
controls the strength of depth emphasis, and a noise standard
deviation $\sigma$ that controls the degree of organic variation.

The three-layer stippling architecture (fine stipple, depth-emphasis,
and edge/background) with cross-layer overlap pruning resolves the
fundamental tension between dot size variation and dot separation,
maintaining maximum blob area below 0.65\% of the canvas while
supporting a dot radius range of 0.5--15.4\,px. Semantic-aware dual-color
rendering, enabled by YOLO26 semantic segmentation, assigns different
dot colors to foreground and background regions, creating a
figure--ground separation that echoes Kusama's compositional practice.

Experiments on three natural images demonstrate that the pipeline achieves
controllable depth encoding, maintains dot separation, and produces
visually compelling results across five high-chroma Kusama palettes.
The ablation study confirms that cross-layer pruning is the most
critical component for preventing blob merging, and that each
layer contributes distinct aesthetic qualities to the final rendering.

This work demonstrates that modern pretrained vision models---originally
developed for autonomous driving and scene understanding---can serve as
artistic primitives for non-photorealistic rendering. The depth and
semantic information they provide opens new dimensions of control
beyond the traditional luminance-based paradigm, enabling computational
approaches to art styles that have previously resisted algorithmic
reproduction. We hope this work encourages further exploration at the
intersection of computational aesthetics, depth estimation, and
artistic style rendering.
